\section{Vision as a Perspective: la visión no es un campo menor}

La sección anterior vinculó la trayectoria de investigación personal con la representation; esta sección explica con claridad «por qué seguir trabajando en vision». Lo esencial aquí no es defender una disciplina antigua, sino redefinir el lugar de la vision en la siguiente generación de inteligencia.

Cuando la entrevista entra en «el mundo de las representaciones», Xie Saining (谢赛宁) plantea una distinción importante: la computer vision no es solo un conjunto de tareas como classification, detection y segmentation, sino una perspective desde la cual mirar la inteligencia. Trabaja con continuous, high-dimensional, noisy signals, es decir, señales en un espacio continuo, de alta dimensión y con ruido. Este tipo de señales es difícil de tokenizar de forma sencilla y tampoco trae de origen etiquetas escritas por personas.

Esto explica por qué el auge de los LLM no lo desanimó. El éxito de los LLM puso en primer plano las interfaces de lenguaje, los sistemas multimodales y el entrenamiento a mayor escala, y con ello le dio a la vision la oportunidad de dejar atrás las tareas aisladas y abordar los problemas amplios de la «inteligencia del mundo real». El verdadero peligro no es que los LLM sean demasiado potentes, sino que todos los problemas de visión se vean obligados a someterse a la narrativa de los modelos de lenguaje y que, al final, la visión quede reducida a un apéndice del prompt y del caption.

\begin{figure}[H]
\centering
\includegraphics[width=0.94\textwidth]{representation-ladder.png}
\caption{De las tareas de visión a los modelos predictivos del mundo: la escalera de capacidades que se desprende de la entrevista.}
\end{figure}

\begin{knowledgebox}{Lectura de la figura: la transición de capacidades de la task al world model}
Los niveles L0 a L4 de la figura no son una clasificación oficial, sino una reconstrucción de la lógica de la entrevista. Las primeras tareas de visión resuelven clasificación, detección y segmentación; la etapa multimodal convierte el lenguaje en la interfaz; más adelante, el sistema debe comprender flujos continuos de eventos, estructura espacial y las consecuencias de cada action. El punto de llegada no es «saber responder preguntas sobre una imagen», sino poder comprimir el flujo visual en una representación del mundo útil para predecir y planificar.
\end{knowledgebox}

\subsection{Por qué importan las representaciones jerárquicas}

Xie Saining menciona en repetidas ocasiones la hierarchical representation. La intuición es la siguiente: un agente no puede memorizar de forma explícita cada píxel, cada molécula y cada parámetro físico del mundo; tiene que aprender a abstraer. Abstraer no significa perder información, sino conservar la información relacionada con la tarea actual, con las acciones futuras y con el costo de las decisiones.

Por ejemplo, la mesa, el micrófono, la iluminación, el sonido y las texturas de una habitación pueden modelarse con todo detalle; pero si el objetivo es continuar la conversación, el sistema solo necesita conocer estados relevantes para la tarea: que el micrófono puede colocarse sobre la mesa, la posición relativa de las dos personas y si el sonido puede captarse. Este state «suficiente, no exhaustivo» es justamente el punto donde se encuentran el representation learning y el world model.

\begin{importantbox}{El problema central del aprendizaje de representaciones}
Una buena representation no reconstruye todos los detalles, sino que comprime señales de alta dimensión en un estado suficiente para sostener la predicción, la planificación y la acción. Debe estar más cerca del mundo físico que las etiquetas de lenguaje y, a la vez, ser más abstracta que los píxeles en bruto.
\end{importantbox}

\subsection{La ayuda y la contaminación del lenguaje}

El lenguaje como interface es sumamente útil. Permite que los sistemas multimodales definan problemas, formulen preguntas y den respuestas en lenguaje natural, y facilita que los sistemas de visión se alineen con los objetivos humanos. Pero el lenguaje también puede convertirse en un shortcut. En la entrevista, Xie Saining describe la tentación que el lenguaje ejerce sobre los sistemas de visión con las imágenes de una «muleta» y del «opio»: el lenguaje hace que el sistema parezca más inteligente, pero puede impedir que desarrolle la capacidad de procesar de verdad un mundo continuo.

\begin{warningbox}{Multimodal no equivale a comprensión real}
Si un benchmark de visión puede responderse correctamente apoyándose sobre todo en el sentido común lingüístico, entonces no demuestra que el modelo comprenda la imagen, el video o el mundo físico. Una vez que interviene el lenguaje, la evaluación debe verificar si el modelo realmente usó representaciones visuales o si resolvió la tarea solo con conocimiento previo textual.
\end{warningbox}

\subsection{Resumen de la sección}

La idea central de Vision as a perspective es esta: la visión no es una tarea menor absorbida por los LLM, sino el problema cognitivo de base que exige la construcción de modelos del mundo. Requiere que el modelo procese señales continuas, abstracciones jerárquicas, estructura espacial, flujos de eventos y las consecuencias reales de las acciones.
